\section{The elementary one-step projection estimate}
\label{sec:leaf-transfer}
We connect the local splitting to the elementary insertion estimate used
in Section~\ref{sec:ope-local-projection}. The analytic input is now the
localized expansion of a real inserted field under one fixed nonlinear
block integral. We state the part of that expansion used here and prove
the remaining conversion and its constants.

Let $\mathcal T$ be the fine-to-coarse density transfer with block ratio
$L$, and let $f$ be one real elementary insertion with anchored fine
support $X$, majorant $c$, and mass $m=m_{\rm fine}(X,c)\ge0$.
Use the same fixed coarse background, threshold and activity alphabet
as the unmarked transition. A localized expansion has real regular
terms $h_i$ on anchored coarse supports $X_i$ and a real flagged
source array $k$ such that
\begin{equation}\label{eq:leaf-raw-identity}
 \mathcal T f=\rho_{\rm coarse}\sum_i\chi_{X_i}h_i
                    +e^{E_{\rm coarse}}\mathcal G^{[1]}(k).
\end{equation}
All regular and marked series must be absolutely dominated, and the
original insertion must be integrable under the fine prior and every
transition slice. These conditions are part of the analytic expansion.
They justify the displayed identity as an identity of actual integrals.
In addition, require the marked identity: the raw forced profile of $k$
is the transferred singleton profile of the insertion in its fixed ordinary
or flagged input slot, with the same background and activity alphabet.
This profile identity is a separate part of the expansion hypothesis.

Split the localized terms into the distinguished principal term, when
present, and the remaining terms. The required bounds have the following
form, in the coarse derivative norm and reserved support mass:
\begin{align}
 \sum_{i\ {
m nonprincipal}}m_+(X_i,\|h_i\|)&\le M,
       \label{eq:leaf-tail-input}\\
 m_+(X_*,\|h_*\|)&\le V\,m/g,
       \label{eq:leaf-principal-input}\\
 \operatorname{mass}(k)&\le D s.
       \label{eq:leaf-flag-input}
\end{align}
Here $g>0$ is the normalization gain, $V\ge0$ a fixed-order derivative
price, $D\ge0$ the insertion debt, and $s\ge0$ the flagged-response
price. For a flagged elementary insertion there is no principal term.
For a regular elementary insertion there is just one, at the fixed
principal location. The indices and supports are selected before
forming complex linear combinations.

\subsection{The principal-term reserve}
The principal estimate \eqref{eq:leaf-principal-input} follows from
three more primitive bounds used by the local integral expansion.
Write $D_X=C_B(X_*)$, let $\omega\ge0$ be its localized majorant,
and suppose
\[
 \|h_*\|\le D\omega\,\Phi(X_*),\qquad
 \omega e^{(a+e)D_X}\le1,\qquad
 D\le m/g,
\]
where $e\ge0$ and $a\ge u+3$. Its derivative price has the form
$\Phi(X_*)\le V e^{D_X/2}$. The latter is the polynomial-to-exponential
bound on the support size: for a fixed integer $r\ge0$ one can take
$C_r=\sup_{z\ge0}z^re^{-z/2}<\infty$ and bound $z^r$ by
$C_re^{z/2}$. The additional fixed factor in $V$ is the spatial
weight of the derivative directions at the given coupling.
Consequently
\[
 e^{(u+2)D_X}\|h_*\|
 \le DV\omega e^{(u+5/2)D_X}
 \le DV\omega e^{(a+e)D_X}
 \le Vm/g.
\]
If $D_X\ge B$, the same calculation leaves a factor $e^{-B/2}$,
because $a+e-(u+5/2)\ge1/2$. This is the stronger principal
estimate used when its support is large. The local derivative
estimate and the weighted bound on $\omega$ remain substantive
inputs from the nonlinear expansion.

\subsection{Complex fields and the raw expansion}
To treat a complex insertion, take two real elementary insertions
with the same support, type and majorant, and combine them as
$f=f_1+if_2$. Their masses are the same $m$; their principal support
and index agree. For $z\in\mathbb C$, define
\[
      H_i=z(h_i^{(1)}+ih_i^{(2)}),\qquad
      K=z(k^{(1)}+ik^{(2)}).
\]
The sum of the two real majorants times $|z|$ dominates each
complex term. Thus the nonprincipal reserved mass is at most
$2|z|M$, the principal reserved mass is at most $2|z|Vm/g$, and
$\operatorname{mass}(K)\le2|z|Ds$.
The raw seed is zero. Pad the common countable list by zero terms
on the singleton anchor support if a natural-number indexing is
needed; neither the represented density nor its cost changes.

Linearity of masks, marked-gas insertion and $\mathcal T$, with
absolute convergence of the series, gives the exact raw density
$z\mathcal T(f_1+if_2)$. This is the reason the same supports and
principal indices were retained. Applying the splitting of
Section~\ref{sec:local-splitting} to these raw terms preserves the
identity. In the elementary application take $z=L^d$ and
$f_2=0$ with the same majorant as $f_1$. The bound with the factor
two remains valid and avoids making the chosen dictionary depend
on whether one component happens to vanish.

\subsection{The projection allowance}
Let $F_d$ be the support-join and jet-extraction factor from
\eqref{eq:local-join-price}, let $C_{\rm orb}\ge1$ be the finite
coefficient-action price, and let $C_{\rm turn}\ge1$ be the price
for transporting the local remainder norm under the required
lattice motions. Put
\[
 H_d=1+2F_d(1+C_{\rm orb}+C_{\rm turn}),\qquad
                     \eta_0=\operatorname{rate}(L)/H_d.
\]
Assume the already chosen coupling satisfies
\begin{equation}\label{eq:leaf-small-inputs}
 2L^d M\le\eta_0m,\qquad
 2\kappa^{-1}L^dDs\le\eta_0m,\qquad
                      (2B+1)^4\kappa\le1,
\end{equation}
where $\kappa=\sqrt{\text{coarse forced allowance}}>0$.
The first two inequalities are required uniformly over insertions.
They are paid by the same coupling delay as the chosen nonlinear
transition; choosing a different history here is not allowed.

Set
\[
         E=2L^dV/g,\qquad e_0=(\eta_0+|E|)m.
\]
Then $e_0\ge\eta_0m\ge0$. The principal bill is at most $e_0$,
and the nonprincipal bill is at most $\eta_0m$. The total raw regular
reserved mass is therefore at most $\eta_0m+e_0\le2e_0$.
Write $p_0=F_d(\eta_0m+e_0)\le2F_de_0$ and
$b_0=2\kappa^{-1}L^dDs\le e_0$.
The local splitting gives a coefficient change at most $p_0$,
a regular remainder of mass at most $p_0$, a normalized forced
remainder of mass at most $b_0+p_0$, and a response of reserved
mass at most $p_0$. To bound the coefficient change by $p_0$,
use $R_d\le F_d$, which follows from $w\ge0$ and the definition
of $F_d$.

It follows directly that
\[
 C_{\rm orb}\mu(q)\le C_{\rm orb}p_0\le H_de_0,
 \qquad
 C_{\rm turn}p_0+(b_0+p_0)\le H_de_0.
\]
The response cost is also at most $H_de_0$. Therefore any
$Q_e$ satisfying
\begin{equation}\label{eq:leaf-projection-allowance}
                  Q_e\ge(\eta_0+|2L^dV/g|)H_d
\end{equation}
gives a coefficient field with coefficient and remainder bounds
$Q_em$ and exact density $L^d\mathcal T f$.
This is the numerical projection bound in
\eqref{eq:ope-leaf-projection}.

The marked transition is retained as well. The original real
insertion occupies its fixed ordinary or flagged slot with cost
$m$; the real/imaginary convention just used gives the available
bound $2m$. Its raw output forced profile is the transferred
singleton profile. Splitting adds exactly the hook of the response
$\psi$ constructed in Section~\ref{sec:local-splitting}, of cost
at most $Q_em$. Thus the transition identity has fine cost $2m$,
coarse response cost $Q_em$, and multiplier $L^d$, precisely as
required in Section~\ref{sec:ope-local-projection}.
This last assertion uses the marked profile hypothesis stated alongside
\eqref{eq:leaf-raw-identity}; the scalar density identity alone would
not suffice.

The analytic work remaining at this point is explicit:
construct the common localized expansion with its singleton identity,
prove \eqref{eq:leaf-tail-input} and \eqref{eq:leaf-flag-input}, and
fund \eqref{eq:leaf-small-inputs} on the selected history.
The finite projection, support correction, principal reserve
conversion, and passage to the elementary allowance have now been
proved from those inputs.
